% SYNC-ID: threats-to-validity
\section{Threats to Validity}\label{sec:threats-to-validity}
% CONTENT_DEFERRED_TO_WRITING_PHASE
\textit{(THIS IS PLACEHOLDER TEXT, PLZ FULFILL IT WHEN WRITING.)}
% SYNC-ID: threats-placeholder-1
Validity limitations will be recorded from source-verified material in the writing phase. This fixture does not assert a threat or a mitigation.
% SYNC-ID: threats-placeholder-2
Each future limitation will be paired with an evidence pointer and a stated scope.
% SYNC-ID: threats-placeholder-3
Separate internal and external validity subsections may be placed here if the approved plan needs them.
% SYNC-ID: threats-placeholder-4
Mitigation language will remain proportional to the evidence available at review time.
% SYNC-ID: threats-placeholder-5
This block reserves a note about replication and transferability without making a prediction.
% SYNC-ID: threats-placeholder-6
Unresolved risks will be listed only after the authors have checked their wording and provenance.
% SYNC-ID: threats-placeholder-7
This section reserves a closing sentence that connects each recorded limitation to the interpretation it qualifies. The sentence will be added only when the corresponding evidence pointer is available.
% SYNC-ID: threats-placeholder-8
The final layout slot is reserved for distinguishing a documented limitation from an unresolved writing task. That distinction will be checked against the same source record as the companion text.
% SYNC-ID: threats-placeholder-8b
The validity section can later use this space to state whether a limitation affects interpretation, transfer, or reproducibility. The finished sentence must identify its boundary without suggesting a remedy that has not been checked. The present text is deliberately noncommittal and exists only to keep the two-column page balance natural.
